% from ega2-1-fr.tex
\oldpage[II]{137}

$x\in X_g\subset f^{-1}(U)$. Alors il existe deux entiers $m>0$, $n>0$, un
homomorphisme
$u:\mathcal{O}_Y^m\to f_*(\mathcal{L}^{\otimes n})$ et un voisinage ouvert
$V$ de $y$ tels que la restriction $u|V$ soit un isomorphisme de
$\mathcal{O}_Y^m|V$ sur $f_*(\mathcal{L}^{\otimes n})|V$.
\end{lemma}

Soient $C$ l'anneau (intègre) de $U$, $k=\mathcal{O}_y$ son corps des
fractions~: comme $f$ est fini, $U'=f^{-1}(U)$ est affine
(\hyperref[II.1.3.2-fr]{1.3.2})~; soit $D$ son anneau (intègre) de corps des
fractions $K=\mathcal{O}_x$~; par hypothèse $D$ est un $C$-module de type
fini (\hyperref[II.6.1.4-fr]{6.1.4}), donc $K$ est une extension de rang fini
de $k$. La fibre
\[
f^{-1}(y)=X\times_Y\operatorname{Spec}(k(y))
=X\times_Y\operatorname{Spec}(\mathcal{O}_y)
\]
s'identifie à $\operatorname{Spec}(K)$
(\hyperref[I.3.6.5-fr]{I, 3.6.6})~; soient $s_i$ ($1\leq i\leq m$) des
éléments de $D$ formant une base de $K$ sur $k$. Il existe $n>0$ tel que les
sections $(s_i|X_g)g^{\otimes n}$ de $\mathcal{L}^{\otimes n}$ au-dessus de
$X_g$ se prolongent en sections $b_i$ ($1\leq i\leq m$) de
$\mathcal{L}^{\otimes n}$ au-dessus de $X$
(\hyperref[I.9.3.1-fr]{I, 9.3.1}). Les $b_i$ sont aussi par définition des
sections de $f_*(\mathcal{L}^{\otimes n})$ au-dessus de $Y$ et définissent
donc un homomorphisme
$u:\mathcal{O}_Y^m\to f_*(\mathcal{L}^{\otimes n})$
(\hyperref[0.5.1.1-fr]{0, 5.1.1})~; nous allons voir que $u$ répond à la
question. On a $\mathcal{L}^{\otimes n}|U'=\widetilde{M}$, où $M$ est un
$D$-module de type fini, donc si $\varphi$ est l'injection $C\to D$
correspondant au morphisme $f^{-1}(U)\to U$ restriction de $f$,
$M_{[\varphi]}$ est un $C$-module de type fini~; par suite
\[
f_*(\mathcal{L}^{\otimes n})|U=(M_{[\varphi]})^\sim
\]
(\hyperref[I.1.6.3-fr]{I, 1.6.3}) est cohérent et comme $U$ est un ouvert
affine quelconque de $Y$, $f_*(\mathcal{L}^{\otimes n})$ est cohérent~; en
outre $u|U=\widetilde{\theta}$, où $\theta$ est un $C$-homomorphisme
$C^m\to M_{[\varphi]}$, et $u_y=\theta_y$ est l'homomorphisme
\[
\theta\otimes1:K^m=C^m\otimes K\longrightarrow
M_{[\varphi]}\otimes K,
\]
or ce dernier est par définition un \emph{isomorphisme}, car les $(b_i)_x$
forment une base de $M_{[\varphi]}\otimes K$ sur $k$, $g_x$ étant par
hypothèse un générateur de $(\mathcal{L}^{\otimes n})_x$. On en conclut que
les supports des $\mathcal{O}_Y$-Modules $\operatorname{Ker}u$ et
$\operatorname{Coker}u$ ne contiennent pas $y$~; comme ces
$\mathcal{O}_Y$-Modules sont cohérents
(\hyperref[0.5.3.3-fr]{0, 5.3.3}), leurs supports sont fermés
(\hyperref[0.5.2.2-fr]{0, 5.2.2}), d'où le lemme.

Cela étant, les hypothèses du lemme
(\hyperref[II.6.7.1.1-fr]{6.7.1.1}) sont remplies dans le cas que nous
considérons en prenant $\mathcal{L}=\mathcal{O}_X$, puisque $X$ est affine
(\hyperref[I.1.1.10-fr]{I, 1.1.10})~; nous poserons
$\mathcal{A}=\mathcal{O}_Y$, $\mathcal{B}=f_*(\mathcal{O}_X)$. En vertu du
critère de Serre (\hyperref[II.5.2.1.1-fr]{5.2.1.1}), il suffit de prouver
que pour tout $\mathcal{O}_Y$-Module cohérent $\mathcal{F}$, on a
$\mathrm{H}^1(Y,\mathcal{F})=0$~; il suffit même de le prouver lorsque
$\mathcal{F}\subset\mathcal{O}_Y$, ce qui entraîne que $\mathcal{F}$ est
sans torsion puisque $Y$ est intègre~; en fait nous allons montrer que
$\mathrm{H}^1(Y,\mathcal{F})=0$ pour \emph{tout} $\mathcal{O}_Y$-Module
cohérent \emph{sans torsion} $\mathcal{F}$. Or, l'homomorphisme
$u:\mathcal{A}^m\to\mathcal{B}$ définit un homomorphisme
\[
v:\mathcal{G}
=\mathcal{H}om_{\mathcal{A}}(\mathcal{B},\mathcal{F})
\longrightarrow
\mathcal{H}om_{\mathcal{A}}(\mathcal{A}^m,\mathcal{F})
=\mathcal{F}^m.
\]
Montrons d'abord que $v$ est \emph{injectif}~: par hypothèse,
$\mathcal{T}=\operatorname{Coker}u$ a un support ne rencontrant pas $V$, donc
est un $\mathcal{O}_Y$-Module de torsion
(\hyperref[I.7.4.6-fr]{I, 7.4.6})~; la suite exacte
\[
\mathcal{A}^m\longrightarrow\mathcal{B}\longrightarrow\mathcal{T}
\longrightarrow0
\]
donne, par exactitude à gauche du foncteur
$\mathcal{H}om_{\mathcal{A}}$, la suite exacte
\[
0\longrightarrow
\mathcal{H}om_{\mathcal{A}}(\mathcal{T},\mathcal{F})
\longrightarrow\mathcal{G}\xrightarrow{v}\mathcal{F}^m.
\]
Mais comme $\mathcal{F}$ est sans torsion, on a
$\mathcal{H}om_{\mathcal{A}}(\mathcal{T},\mathcal{F})=0$
(\hyperref[0.5.2.6-fr]{0, 5.2.6}), d'où notre assertion. On a donc la suite
exacte
\[
0\longrightarrow\mathcal{G}\longrightarrow\mathcal{F}^m
\longrightarrow\operatorname{Coker}v\longrightarrow0
\]

